\section{Threats to Validity}

This section follows the four-category taxonomy of Wohlin et al.~\cite{wohlin2012experimentation} for empirical software engineering studies, extended with statistical considerations.

\subsection{Conclusion Validity}

\textbf{Statistical method.} All comparisons use BCa bootstrap confidence intervals~\cite{efron1994introduction} across 25 evaluation seeds with no null hypothesis significance testing. Evidence is based on CI non-overlap, which is a conservative but defensible criterion. A limitation of this approach is that CIs can overlap while a permutation test would reject the null---we err on the side of requiring non-overlap for a claim of substantive evidence.

\textbf{Multiple comparisons.} Results for nine methods at three $K$ values and two metrics (P@$K$, NDCG@$K$) constitute 54 comparisons. No familywise or false discovery rate correction is applied; all comparisons are pre-registered and the paper's primary claims are limited to the largest, most stable differences (HygienePrio-full vs.\ EPSS-only, $\geq$21 pp at all $K$). Smaller differences (interaction term, freshness dimension) are explicitly treated as equivocal.

\textbf{Seed budget.} With 25 evaluation seeds, the bootstrap CIs have limited resolution. Differences of $< \approx$3 pp are unlikely to be distinguishable from seed-to-seed variation (estimated from the within-seed variance of the most stable comparisons). All headline claims involve differences substantially larger than this threshold ($\geq$21 pp). The 5-seed calibration set is thin; weight calibration on 5 seeds may produce noisy estimates relative to the true optimal weights.

\textbf{Bootstrap variance.} The BCa CI computation uses 10,000 bootstrap resamples per comparison. For a 25-seed evaluation, bootstrap variance in CI boundaries is small relative to the observed differences. The BCa procedure corrects for bias and skewness in the bootstrap distribution, making it preferable to the percentile bootstrap for skewed P@$K$ distributions.

\subsection{Internal Validity}

\textbf{Synthetic data generation.} The EEHDA fleet generator uses structural priors (DBIR patch-lag statistics, NVD severity distributions) that represent population-level aggregates. If real fleets deviate substantially from these priors---for example, if patch debt is correlated with host criticality in ways not modeled---results may not replicate.

\textbf{Ground truth circularity.} The pre-registered ground truth includes HRS $>$ 75th percentile as a labelling condition (\S3). This provides a structural advantage to HygienePrio over baselines that do not incorporate HRS. The threshold (75th percentile) was fixed before evaluation, preventing post-hoc optimization, but the choice of this specific threshold is not externally validated. Results under the 60th or 90th percentile threshold are not reported; sensitivity to this choice is an open question.

\textbf{Threshold choices.} The EPSS $>$ 0.10 and HRS $>$ 75th percentile thresholds that define the positive class are operationally motivated but not derived from validated exploitation data. These thresholds directly control the base rate of true positives (6.2\% mean across seeds) and therefore affect P@$K$ values for all methods. The relative ordering of methods may be less sensitive to threshold choice than the absolute P@$K$ values, but this is not verified.

\textbf{Calibration leakage.} Calibration and evaluation seeds are strictly separated (5 calibration, 25 evaluation, from adjacent seed values). If the generator has strong auto-correlation across adjacent seed values, the calibration/evaluation separation may be partially compromised. Generator analysis supporting independence of adjacent seeds is not reported; this is a known limitation.

\textbf{Measurement validity of HRS dimensions.} PatchPosture is measured as a static snapshot of the proportion of applicable CVEs that remain unpatched. This does not capture patch scheduling latency: a host may have 50\% unpatched CVEs today with all patches scheduled for tomorrow. In real environments, operational scheduling context would be needed to distinguish genuine neglect from planned maintenance.

\subsection{External Validity}

\textbf{Generalizability to real enterprise fleets.} The synthetic fleet cannot replicate the diversity of real enterprise environments. All results are bounded to the synthetic evaluation context. Claims about real-fleet performance require studies on real fleet data with appropriate institutional approvals and privacy protections.

\textbf{Generalizability across sectors.} The EEHDA generator is calibrated for government-shaped public-sector fleet characteristics (following Paper~1). Commercial enterprise, healthcare, and critical infrastructure fleets may have substantially different hygiene posture distributions, making the specific calibrated weights inapplicable without recalibration.

\textbf{Temporal validity.} EPSS scores change daily. The evaluation uses a fixed snapshot per seed. In a production deployment, HygienePrio would use daily-updated EPSS scores; the evaluation captures only one temporal cross-section per seed, which may not reflect the distribution of daily prioritization decisions over time.

\subsection{Construct Validity}

\textbf{HRS as a proxy for real risk.} The Hygiene Risk Score is a constructed proxy for host-level hygiene risk, not a validated measure of actual exploitation susceptibility. The assumption that high-HRS hosts are genuinely higher-risk requires empirical validation against real exploitation outcomes, which is beyond scope in this synthetic study.

\textbf{Dimension underrepresentation.} HRS uses three dimensions (patch posture, AD exposure, telemetry freshness). Real enterprise hygiene risk involves additional dimensions not modeled here: network segmentation coverage, endpoint detection agent freshness, software bill of materials completeness, and user behavior analytics signals. The three-dimension HRS may systematically underestimate risk for certain host archetypes (e.g., network-isolated hosts with good patch posture but no EDR coverage).

\textbf{P@$K$ as a remediation metric.} Precision@$K$ assumes equally valuable true positives. In real operations, some (host, CVE) pairs may be substantially more consequential than others (e.g., domain controllers vs.\ user workstations), and a severity-weighted metric might yield different method rankings. Oracle-gap captures ceiling relative to the positive base rate but shares the binary relevance limitation.

\textbf{Scorer linearity.} The linear scorer with one interaction term may not capture complex non-linear relationships between hygiene dimensions and actual exploitation risk in real fleets.
